\section{From identified coordinates to sharing decisions}
\label{sec:sharing-dynamics}

Once the factors are known, choosing which layers to tie still requires
an objective. We compare three choices and ask whether their preferred
partitions can diverge during training. Let $L=Km$ and let $\mathcal H_m$ contain binary $L\times K$
assignment matrices with one entry per row and $m$ per column. For
$H\in\mathcal H_m$, $Q_H=HH^\top/m$ projects onto group-constant rows.
Partitions are compared without group labels. Define
\begin{align}
 S_A(H)&=\langle A,H\rangle, &
 J_\Theta(H)&=\|(I-Q_H)AB\|_F^2, &
 J_A(H)&=\|(I-Q_H)A\|_F^2. \label{eq:three-decisions}
\end{align}
\emph{Coefficient assignment} maximizes $S_A$: it assigns to fixed simplex
vertices, since $\|A-H\|_F^2=\|A\|_F^2+L-2S_A(H)$.
\emph{Weight clustering} minimizes $J_\Theta$, refitting each group center;
it is router clustering in the metric $BB^\top$.
\emph{Router clustering} minimizes $J_A$ with free centers. These are
different objectives, even when the original factors are known exactly.

\subsection{What response information is sufficient for a decision?}

After fixing $H$, optimize only the hard-shared bases with rate $\eta_B$.
Their effective response is
$\mathcal R_H=\eta_B(HH^\top)\otimes I_D$; it does not depend on their
initial fitted values. Here $\|\cdot\|_{\rm HS}$ denotes the Frobenius norm
of the full $LD\times LD$ operator, distinct from $\|\cdot\|_{F\to F}$.

\begin{theorem}[Response-optimal balanced tying]
\label{thm:response-partition}
For the soft response in Equation~\eqref{eq:response_blocks} and hard
responses with the same $\eta_B>0$, there is a scalar $c(A,B)$ independent
of $H\in\mathcal H_m$ such that
\begin{equation}
 \|\mathcal R_{A,B}-\mathcal R_H\|_{\rm HS}^2
 =c(A,B)+2\eta_B^2 Dm J_A(H). \label{eq:response-partition}
\end{equation}
Thus the global minimizers are exactly the balanced router-clustering
partitions. Off-diagonal entries of $AA^\top$ suffice to select them.
\end{theorem}
\begin{proof}
Set $\gamma=\eta_Z/\tau^2$ and $T_i=B^\top C_i^2B$.
Blockwise expansion gives
\begin{align*}
 \|\mathcal R_{A,B}-\mathcal R_H\|_{\rm HS}^2
 ={}&\eta_B^2D\|AA^\top-HH^\top\|_F^2\\
 &+2\eta_B\gamma\sum_i(\|a_i\|^2-1)\operatorname{tr}T_i
   +\gamma^2\sum_i\|T_i\|_F^2.
\end{align*}
The last two terms are partition independent. Use
$\|HH^\top\|_F^2=Lm$ and
$\langle AA^\top,HH^\top\rangle=\|H^\top A\|_F^2
=m(\|A\|_F^2-J_A(H))$. Diagonal Gram entries contribute the same sum
for every partition, proving the final assertion.
\end{proof}

The reduction uses standard Gram-based clustering algebra
\citep{dhillon2004kernel}. Under the equal-capacity and learning-rate
assumptions, the full softmax covariance contribution is independent of
the partition. The metric averages
squared response error over isotropic effective gradients. It need not select
the best partition for actual task gradients, operator norm, AdamW, or final
loss. Unlike full factor identification, this decision does not require
recovering $A$: the off-diagonal response blocks already provide the
sufficient Gram entries. Response-Gram clustering is therefore a simpler
sufficient alternative for this objective.

\paragraph{Equivalence and strict conflicts.}
For $K=2$ and $b_1\ne b_2$, coefficient assignment and global balanced
weight clustering have exactly the same optimal unlabeled partitions.
Writing $a_i=(t_i,1-t_i)$, both select the top $L/2$ values of $t_i$ and
their complement. This includes ties and is a global-objective statement.
For arbitrary $K$, if
$P_0BB^\top P_0=\lambda P_0$, where
$P_0=I-\ones\ones^\top/K$ and $\lambda>0$, then
$J_\Theta=\lambda J_A$. This equates the two free-center clustering
objectives, not coefficient assignment.

For $K=3$, strictly interior full-rank examples separate all these claims
(Appendix~\ref{app:decision-proofs}). One example with $L=6,D=9$ and
basis condition number four has a unique weight-optimal partition with
$(J_\Theta,J_A)=(65,50)/529$, whereas the unique response-optimal
partition has $(119,14)/529$. A separate, exact product-preserving change
of factors gives two identical effective models with \emph{different unique}
response-optimal partitions. Consequently no decision using only $\Theta$
can always be response optimal. Response information can therefore matter even when full coordinate
reconstruction is unnecessary.

\subsection{Training reachability and finite-time limits}

The static examples leave open whether training can move from agreement
to disagreement. The following existence result shows that it can.

\begin{theorem}[A realizable transverse decision crossing]
\label{thm:decision-crossing}
There is a fixed squared loss $\ell(\Theta)=\tfrac12\|\Theta-T\|_F^2$
and a nonempty open set of initial softmax/shared-basis states with
$K=3,L=6,D=9$ such that, under Euclidean factor gradient flow with
$\eta_B=\eta_Z=\tau=1$, coefficient assignment and global balanced weight
clustering start with the same unique partition and subsequently disagree.
The weight partition stays fixed through the crossing; both factors remain
full rank and the router stays strictly positive on the considered interval.
\end{theorem}

\paragraph{Mechanism.}
Appendix~\ref{app:decision-crossing-proof} gives rational factors at a tie
between exactly two coefficient assignments, $H_p,H_q$, with a strict
weight-clustering gap. Put $E=H_p-H_q$ and set the fixed loss gradient at
that point to $g_i=B^\top C_i^2 e_i$ by choosing $T=AB-G$. Then
\begin{equation}
 \left.\frac{d}{dt}\langle A,H_p-H_q\rangle\right|_{t=0}
 =-\sum_i\|B^\top C_i^2e_i\|^2<0. \label{eq:crossing-speed}
\end{equation}
Local ODE existence and the strict remaining gaps give a crossing while
the weight optimum stays fixed. Taking a nearby earlier point as the
initial state produces ordinary forward training under the same loss.
Continuous dependence yields an open neighborhood of such initial states.
The constructed target proves that the crossing is reachable. Its
frequency and attraction properties under natural data remain open.

Positive margins can also protect agreement for a finite time.
Appendix~\ref{app:decision-stability} gives explicit bounds using initial
assignment/clustering gaps and bounds on $\|B\|_2$ and
$\|\nabla\ell\|_F$. Under these conditions, agreement persists for the bounded interval;
it need neither hold indefinitely nor inevitably break down.

\paragraph{Why static and dynamic objectives can both matter.}
Consider soft and hard trajectories of the same loss. Assume its gradient
is $L_\ell$-Lipschitz between the trajectories,
$\|\nabla\ell(\Theta_s(t))\|_F\le M$, and
$\|\mathcal R_s(t)-\mathcal R_H\|_{F\to F}\le\delta$ throughout the
interval. With $c=\eta_BmL_\ell$ and group-mean initialization,
\begin{equation}
 \|\Theta_s(t)-\Theta_h(t)\|_F
 \le e^{ct}\sqrt{J_\Theta(H)}
       +M\delta\frac{e^{ct}-1}{c}. \label{eq:trajectory-sharing-bound}
\end{equation}
The last quotient is $t$ at $c=0$. This standard trajectory perturbation
bound, proved in Appendix~\ref{app:decision-stability}, separates initial
weight distortion from continuing response distortion. Matching training
kernels is already a compression objective~\citep{pmlr-v119-liu20o}.
Our instantaneous Hilbert--Schmidt optimum only bounds operator mismatch
at that time; it does not supply a uniform-in-time $\delta$ or guarantee
lower downstream loss.
